% IFJ + IAL projekt 2026/2027
% dokumentace.tex
% Richard Komora (xkomorr00)
% Martin Huleja (xhulejm00)

\documentclass[a4paper, 11pt]{article}
\usepackage[left=2cm, text={17cm, 24cm}, top=3cm]{geometry}
\usepackage[T1]{fontenc}
\usepackage[utf8]{inputenc}
\usepackage[dvips, unicode, hidelinks, bookmarksopen]{hyperref}
\usepackage{xltabular}
\usepackage{pdflscape}
\usepackage{graphicx}
\usepackage{booktabs}

% Command definitions
\newcommand{\tableheader} {
  \hline
  Číslo & $A \rightarrow \alpha$ & $Predict (A \rightarrow \alpha)$ \\
  \hline
}

\newcommand{\rot}[1]{\rotatebox{90}{#1}}


\begin{document}

Členovia tímu:
\textbf{Martin Huleja xhulejm00:282792}
Adam Havier xhaviea00:281928
Richard Komora xkomorr00:281910
Enzo Štětina xstetie00:282488

Varianta: TRP-izp

Rozdelenie bodov:
282792:xhulejm00:25
281928:xhaviea00:25
281910:xkomorr00:25
282488:xstetie00:25

//TODO: identifikátory implementovaných rozšíření,

GitHub Repozitár: \href{https://github.com/martinhuleja/IFJ\_project}{https://github.com/martinhuleja/IFJ\_project}

//TODO: priebežne dopĺňať
Rozdelenie práce:
Martin Huleja - Parser
Adam Havier - Generátor IFJcode26
Richard Komora - Tabuľka symbolov + sémantika
Enzo Štětina - Lexer + automatické spúštanie testov

//TODO:
stručný popis způsobu spolupráce, kultury code review a schvalování PR (max. jednotky vět
či odrážek); stručný popis využitých forem testování (max. jednotky vět či odrážek); zvolený
styl formátování kódu;
• blokové schéma architektury vašeho řešení, které bude obsahovat jednotlivé moduly překla-
dače (např. lexikální, syntaktická, sémantická analýza atp.) a musí zobrazovat i datové toky
mezi jednotlivými komponentami – specifikujte, v jakých strukturách (např. tokeny, AST, me-
zikód) si komponenty předávají informace. Z diagramu a jeho textového popisu musí být jasné
vstupy, výstupy, význam komponent a jejich propojení;
• diagram automatu, který specifikuje lexikální analyzátor (podle konvencí přednášených v IFJ;
s případným slovním popisem odchylek od standardního modelu konečného automatu);
• precedenční tabulku, podle kterých jste imple-
mentovali syntaktický analyzátor;
• stručný popis každého modulu (informačně obohacující věta na dvojici .c/.h);
• bibliografické citace všech převzatých zdrojů (vč. kódu, algoritmů);
• popis využití prostředků AI (viz sekci 4).

\section{LL gramatika}

Pro syntaktickou analýzu shora dolů byla využita metoda rekurzivního sestupu. \\

\noindent $G = (N, T, P, S)$

\begin{itemize}
  \item $N = \{$ Program, Prologue, FunctionList, FunctionListRest, FunctionDef, MainOpt,
          ParamGroup, ParamListOpt, ParamList, ParamListRest, Param, TypeOpt, Type,
          FunctionBody, Block, StatementList, StatementListRest, Statement, StatementEnd,
          VarDef, VarKind, IfStatement, WhileStatement, ReturnStatement, ArgGroup, ArgList,
          ArgListOpt, ArgListRest, Expr, ExprOrCall, ReturnExprOpt, CallOpt, EOLOrEOF,
          StatementAfterId, ExprNonId, ExprTail, EOLOpt $\}$
  \item $T = \{$ Identifier, IntLiteral, DoubleLiteral, StringLiteral, UnitLiteral, TypeInt,
          TypeDouble, TypeString, TypeUnit, Val, Var, If, Then, Else, While, Do, Return,
          AtMain, Def, IFJ26Id, Import, Assign, LParen, RParen, Colon, Comma, Semicolon,
          Dot, Star, Indent, Dedent, EOL, EOF  $\}$
  \item $S = program$
\end{itemize}

\subsection{Přepisovací pravidla}

\begin{xltabular}{\textwidth}{|l|X|p{0.3\textwidth}|}
  \tableheader
  \endfirsthead

  \tableheader
  \endhead

  1   & Program $\rightarrow$ Prologue FunctionList EOF
      & $\{$ Import $\}$ \\ \hline
  2   & Prologue $\rightarrow$ Import IFJ26Id Dot Star EOL
      & $\{$ Import $\}$ \\ \hline
  3   & FunctionList $\rightarrow$ FunctionDef FunctionListRest
      & $\{$ AtMain, Def $\}$ \\ \hline
  4   & FunctionListRest $\rightarrow$ FunctionDef FunctionListRest
      & $\{$ AtMain, Def $\}$ \\ \hline
  5   & FunctionListRest $\rightarrow \varepsilon$
      & $\{$ EOF $\}$ \\ \hline
  6   & FunctionDef $\rightarrow$ MainOpt EOLOpt Def Identifier ParamGroup TypeOpt Assign FunctionBody
      & $\{$ AtMain, Def $\}$ \\ \hline
  7   & MainOpt $\rightarrow$ AtMain
      & $\{$ AtMain $\}$ \\ \hline
  8   & MainOpt $\rightarrow \varepsilon$
      & $\{$ Def $\}$ \\ \hline
  9   & ParamGroup $\rightarrow$ LParen ParamListOpt RParen
      & $\{$ LParen $\}$ \\ \hline
  10  & ParamListOpt $\rightarrow$ ParamList
      & $\{$ Identifier $\}$ \\ \hline
  11  & ParamListOpt $\rightarrow \varepsilon$
      & $\{$ RParen $\}$ \\ \hline
  12  & ParamList $\rightarrow$ Param ParamListRest
      & $\{$ Identifier $\}$ \\ \hline
  13  & ParamListRest $\rightarrow$ Comma Param ParamListRest
      & $\{$ Comma $\}$ \\ \hline
  14  & ParamListRest $\rightarrow \varepsilon$
      & $\{$ RParen $\}$ \\ \hline
  15  & Param $\rightarrow$ Identifier Colon Type
      & $\{$ Identifier $\}$ \\ \hline
  16  & TypeOpt $\rightarrow$ Colon Type
      & $\{$ Colon $\}$ \\ \hline
  17  & TypeOpt $\rightarrow \varepsilon$
      & $\{$ Assign $\}$ \\ \hline
  18  & Type $\rightarrow$ TypeInt
      & $\{$ TypeInt $\}$ \\ \hline
  19  & Type $\rightarrow$ TypeDouble
      & $\{$ TypeDouble $\}$ \\ \hline
  20  & Type $\rightarrow$ TypeString
      & $\{$ TypeString $\}$ \\ \hline
  21  & Type $\rightarrow$ TypeUnit
      & $\{$ TypeUnit $\}$ \\ \hline
  22  & FunctionBody $\rightarrow$ ExprOrCall EOLOrEOF
      & $\{$ Identifier, IntLiteral, DoubleLiteral, StringLiteral, UnitLiteral, LParen $\}$ \\ \hline
  23  & EOLOrEOF $\rightarrow$ EOL
      & $\{$ EOL $\}$ \\ \hline
  24  & EOLOrEOF $\rightarrow$ EOF
      & $\{$ EOF $\}$ \\ \hline
  25  & FunctionBody $\rightarrow$ EOL Block
      & $\{$ EOL $\}$ \\ \hline
  26  & Block $\rightarrow$ Indent StatementList Dedent
      & $\{$ Indent $\}$ \\ \hline
  27  & StatementList $\rightarrow$ Statement StatementEnd StatementListRest
      & $\{$ Val, Var, If, While, Return, Identifier, IntLiteral, DoubleLiteral, StringLiteral, UnitLiteral, LParen $\}$ \\ \hline
  28  & StatementListRest $\rightarrow$ Statement StatementEnd StatementListRest
      & $\{$ Val, Var, If, While, Return, Identifier, IntLiteral, DoubleLiteral, StringLiteral, UnitLiteral, LParen $\}$ \\ \hline
  29  & StatementListRest $\rightarrow \varepsilon$
      & $\{$ Dedent $\}$ \\ \hline
  30  & StatementEnd $\rightarrow$ EOL
      & $\{$ EOL $\}$ \\ \hline
  31  & StatementEnd $\rightarrow$ Semicolon
      & $\{$ Semicolon $\}$ \\ \hline
  32  & Statement $\rightarrow$ VarDef
      & $\{$ Val, Var $\}$ \\ \hline
  33  & Statement $\rightarrow$ IfStatement
      & $\{$ If $\}$ \\ \hline
  34  & Statement $\rightarrow$ WhileStatement
      & $\{$ While $\}$ \\ \hline
  35  & Statement $\rightarrow$ ReturnStatement
      & $\{$ Return $\}$ \\ \hline
  36  & Statement $\rightarrow$ Identifier StatementAfterId
      & $\{$ Identifier $\}$ \\ \hline
  37  & Statement $\rightarrow$ ExprNonId
      & $\{$ IntLiteral, DoubleLiteral, StringLiteral, UnitLiteral, LParen $\}$ \\ \hline
  38  & StatementAfterId $\rightarrow$ Assign ExprOrCall
      & $\{$ Assign $\}$ \\ \hline
  39  & StatementAfterId $\rightarrow$ ArgGroup
      & $\{$ LParen $\}$ \\ \hline
  40  & StatementAfterId $\rightarrow$ ExprTail
      & $\{$ EOL, Semicolon $\}$ \\ \hline
  41  & ExprNonId $\rightarrow$ IntLiteral ExprTail
      & $\{$ IntLiteral $\}$ \\ \hline
  42  & ExprNonId $\rightarrow$ DoubleLiteral ExprTail
      & $\{$ DoubleLiteral $\}$ \\ \hline
  43  & ExprNonId $\rightarrow$ StringLiteral ExprTail
      & $\{$ StringLiteral $\}$ \\ \hline
  44  & ExprNonId $\rightarrow$ UnitLiteral ExprTail
      & $\{$ UnitLiteral $\}$ \\ \hline
  45  & ExprNonId $\rightarrow$ LParen Expr RParen ExprTail
      & $\{$ LParen $\}$ \\ \hline
  46  & VarDef $\rightarrow$ VarKind Identifier TypeOpt Assign ExprOrCall
      & $\{$ Val, Var $\}$ \\ \hline
  47  & VarKind $\rightarrow$ Val
      & $\{$ Val $\}$ \\ \hline
  48  & VarKind $\rightarrow$ Var
      & $\{$ Var $\}$ \\ \hline
  49  & IfStatement $\rightarrow$ If Expr Then EOL Block EOL Else EOL Block
      & $\{$ If $\}$ \\ \hline
  50  & WhileStatement $\rightarrow$ While Expr Do EOL Block
      & $\{$ While $\}$ \\ \hline
  51  & ReturnStatement $\rightarrow$ Return ReturnExprOpt
      & $\{$ Return $\}$ \\ \hline
  52  & ReturnExprOpt $\rightarrow$ ExprOrCall
      & $\{$ Identifier, IntLiteral, DoubleLiteral, StringLiteral, UnitLiteral, LParen $\}$ \\ \hline
  53  & ReturnExprOpt $\rightarrow \varepsilon$
      & $\{$ EOL, Semicolon $\}$ \\ \hline
  54  & ExprOrCall $\rightarrow$ Identifier CallOpt
      & $\{$ Identifier $\}$ \\ \hline
  55  & ExprOrCall $\rightarrow$ Expr
      & $\{$ IntLiteral, DoubleLiteral, StringLiteral, UnitLiteral, LParen $\}$ \\ \hline
  56  & CallOpt $\rightarrow$ ArgGroup
      & $\{$ LParen $\}$ \\ \hline
  57  & CallOpt $\rightarrow \varepsilon$
      & $\{$ EOL, EOF, Semicolon $\}$ \\ \hline
  58  & ArgGroup $\rightarrow$ LParen ArgListOpt RParen
      & $\{$ LParen $\}$ \\ \hline
  59  & ArgListOpt $\rightarrow$ ArgList
      & $\{$ Identifier, IntLiteral, DoubleLiteral, StringLiteral, UnitLiteral, LParen $\}$ \\ \hline
  60  & ArgListOpt $\rightarrow \varepsilon$
      & $\{$ RParen $\}$ \\ \hline
  61  & ArgList $\rightarrow$ Expr ArgListRest
      & $\{$ Identifier, IntLiteral, DoubleLiteral, StringLiteral, UnitLiteral, LParen $\}$ \\ \hline
  62  & ArgListRest $\rightarrow$ Comma Expr ArgListRest
      & $\{$ Comma $\}$ \\ \hline
  63  & ArgListRest $\rightarrow \varepsilon$
      & $\{$ RParen $\}$ \\ \hline
  64  & ExprTail $\rightarrow \varepsilon$
      & $\{$ EOL, Semicolon $\}$ \\ \hline
  65  & EOLOpt $\rightarrow$ EOL
      & $\{$ EOL $\}$ \\ \hline
  66  & EOLOpt $\rightarrow \varepsilon$
      & $\{$ Def $\}$ \\ \hline

\end{xltabular}

Neterminál Expr je vyhodnocován samostatnou precedenční syntaktickou analýzou zdola nahoru podle precedenční tabulky.

% ====================
% ===== LL table =====
% ====================

\subsection{LL tabulka}

\begin{landscape}
  \begin{table}[p]
  \centering
  \small
  \setlength{\tabcolsep}{4pt}

  \begin{tabular}{l *{33}{c}}
  \toprule

  % === HEADER ===
  \textbf{Neterminál} & \rot{Identifier} & \rot{IntLiteral} & \rot{DoubleLiteral} &
  \rot{StringLiteral} & \rot{UnitLiteral} & \rot{TypeInt} & \rot{TypeDouble} &
  \rot{TypeString} & \rot{TypeUnit} & \rot{Val} & \rot{Var} & \rot{If} & \rot{Then} &
  \rot{Else} & \rot{While} & \rot{Do} & \rot{Return} & \rot{AtMain} & \rot{Def} &
  \rot{IFJ26Id} & \rot{Import} & \rot{Assign} & \rot{LParen} & \rot{RParen} &
  \rot{Colon} & \rot{Comma} & \rot{Semicolon} & \rot{Dot} & \rot{Star} &
  \rot{Indent} & \rot{Dedent} & \rot{EOL} & \rot{EOF} \\
  \midrule
  % === TABLE CONTENT ===
  Program & & & & & & & & & & & & & & & & & & & & & 1 & & & & & & & & & & & & \\
  Prologue & & & & & & & & & & & & & & & & & & & & & 2 & & & & & & & & & & & & \\
  FunctionList & & & & & & & & & & & & & & & & & & 3 & 3 & & & & & & & & & & & & & & \\
  FunctionListRest & & & & & & & & & & & & & & & & & & 4 & 4 & & & & & & & & & & & & & & 5 \\
  FunctionDef & & & & & & & & & & & & & & & & & & 6 & 6 & & & & & & & & & & & & & & \\
  MainOpt & & & & & & & & & & & & & & & & & & 7 & 8 & & & & & & & & & & & & & & \\
  ParamGroup & & & & & & & & & & & & & & & & & & & & & & & 9 & & & & & & & & & & \\
  ParamListOpt & 10 & & & & & & & & & & & & & & & & & & & & & & & 11 & & & & & & & & & \\
  ParamList & 12 & & & & & & & & & & & & & & & & & & & & & & & & & & & & & & & & \\
  ParamListRest & & & & & & & & & & & & & & & & & & & & & & & & 14 & & 13 & & & & & & & \\
  Param & 15 & & & & & & & & & & & & & & & & & & & & & & & & & & & & & & & & \\
  TypeOpt & & & & & & & & & & & & & & & & & & & & & & 17 & & & 16 & & & & & & & & \\
  Type & & & & & & 18 & 19 & 20 & 21 & & & & & & & & & & & & & & & & & & & & & & & & \\
  FunctionBody & 22 & 22 & 22 & 22 & 22 & & & & & & & & & & & & & & & & & & 22 & & & & & & & 25 & & & \\
  Block & & & & & & & & & & & & & & & & & & & & & & & & & & & & & & 26 & & & \\
  StatementList & 27 & 27 & 27 & 27 & 27 & & & & & 27 & 27 & 27 & & & 27 & & 27 & & & & & & 27 & & & & & & & & & & \\
  StatementListRest & 28 & 28 & 28 & 28 & 28 & & & & & 28 & 28 & 28 & & & 28 & & 28 & & & & & & 28 & & & & & & & & 29 & & \\
  Statement & 36 & 37 & 37 & 37 & 37 & & & & & 32 & 32 & 33 & & & 34 & & 35 & & & & & & 37 & & & & & & & & & & \\
  StatementEnd & & & & & & & & & & & & & & & & & & & & & & & & & & & 31 & & & & & 30 & \\
  VarDef & & & & & & & & & & 46 & 46 & & & & & & & & & & & & & & & & & & & & & & \\
  VarKind & & & & & & & & & & 47 & 48 & & & & & & & & & & & & & & & & & & & & & & \\
  IfStatement & & & & & & & & & & & & 49 & & & & & & & & & & & & & & & & & & & & & \\
  WhileStatement & & & & & & & & & & & & & & & 50 & & & & & & & & & & & & & & & & & & \\
  ReturnStatement & & & & & & & & & & & & & & & & & 51 & & & & & & & & & & & & & & & & \\
  ArgGroup & & & & & & & & & & & & & & & & & & & & & & & 58 & & & & & & & & & & \\
  ArgList & 61 & 61 & 61 & 61 & 61 & & & & & & & & & & & & & & & & & & 61 & & & & & & & & & & \\
  ArgListOpt & 59 & 59 & 59 & 59 & 59 & & & & & & & & & & & & & & & & & & 59 & 60 & & & & & & & & & \\
  ArgListRest & & & & & & & & & & & & & & & & & & & & & & & & 63 & & 62 & & & & & & & \\
  Expr & & & & & & & & & & & & & & & & & & & & & & & & & & & & & & & & & \\
  ExprOrCall & 54 & 55 & 55 & 55 & 55 & & & & & & & & & & & & & & & & & & 55 & & & & & & & & & & \\
  ReturnExprOpt & 52 & 52 & 52 & 52 & 52 & & & & & & & & & & & & & & & & & & 52 & & & & 53 & & & & & 53 & \\
  CallOpt & & & & & & & & & & & & & & & & & & & & & & & 56 & & & & 57 & & & & & 57 & 57 \\
  EOLOrEOF & & & & & & & & & & & & & & & & & & & & & & & & & & & & & & & & 23 & 24 \\
  StatementAfterId & & & & & & & & & & & & & & & & & & & & & & 38 & 39 & & & & 40 & & & & & 40 & \\
  ExprNonId & & 41 & 42 & 43 & 44 & & & & & & & & & & & & & & & & & & 45 & & & & & & & & & & \\
  ExprTail & & & & & & & & & & & & & & & & & & & & & & & & & & & 64 & & & & & 64 & \\
  EOLOpt & & & & & & & & & & & & & & & & & & & 66 & & & & & & & & & & & & & 65 & \\
  \bottomrule
  \end{tabular}
  \end{table}
\end{landscape}

\section{Precedenční syntaktická analýza výrazů}

Pro vyhodnocování výrazů byla v souladu se zadáním využita precedenční syntaktická analýza.

\subsection{Operátorová gramatika výrazů}

Gramatika zpracovávající výrazy je definována jako
$G_{\textnormal{expr}} = (\{E\}, T_{\textnormal{prec}}, P_{\textnormal{prec}}, E)$,
kde $T_{\textnormal{prec}} = \{ i, -_{\textnormal{un}}, *, /, +, -_{\textnormal{bi}},
<, <=, >, >=, ==, !=, (, ), \$ \}$

\begin{itemize}
    \item $i$ značí identifikátory a literály typů \texttt{Int}, \texttt{Double}, \texttt{String}, \texttt{Unit},
    \item $-_{\textnormal{un}}$ značí unární operátor minus, $-_{\textnormal{bi}}$ binární odčítání,
    \item $\$$ značí ukončovač výrazu (\texttt{EOL}, \texttt{EOF}, \texttt{then}, \texttt{do}, \dots).
\end{itemize}

\noindent Množina pravidel $P_{\textnormal{prec}}$ zahrnuje: \quad $E \rightarrow i, \qquad E \rightarrow (E),
\qquad E \rightarrow -E, \qquad E \rightarrow E * E, \qquad E \rightarrow E / E, \qquad E \rightarrow E + E,
\qquad E \rightarrow E - E, \qquad E \rightarrow E \ rel \ E$, kde $rel \in \{ <, <=, >, >=, ==, != \}.$

\subsection{Precedenční tabulka}
Hierarchie operátorů vychází ze specifikace jazyka IFJ26.

\begin{table}[htbp]
  \centering
  \begin{tabular}{l *{9}{|c}}

  \toprule
  \textbf{Vrchol / Vstup}
  & $i$ & $-_{\textnormal{un}}$ & $*$, $/$ & $+$, $-_{\textnormal{bi}}$
  & $<, <=, >, >=$ & $==, !=$ & $($ & $)$ & $\$$ \\

  \midrule

  $i$                        &     &     & $>$ & $>$ & $>$ & $>$ &     & $>$ & $>$ \\
  $-_{\textnormal{un}}$      & $<$ & $<$ & $>$ & $>$ & $>$ & $>$ & $<$ & $>$ & $>$ \\
  $*$, $/$                   & $<$ & $<$ & $>$ & $>$ & $>$ & $>$ & $<$ & $>$ & $>$ \\
  $+$, $-_{\textnormal{bi}}$ & $<$ & $<$ & $<$ & $>$ & $>$ & $>$ & $<$ & $>$ & $>$ \\
  $<, <=, >, >=$             & $<$ & $<$ & $<$ & $<$ &     & $>$ & $<$ & $>$ & $>$ \\
  $==, !=$                   & $<$ & $<$ & $<$ & $<$ & $<$ &     & $<$ & $>$ & $>$ \\
  $($                        & $<$ & $<$ & $<$ & $<$ & $<$ & $<$ & $<$ & $=$ &     \\
  $)$                        &     &     & $>$ & $>$ & $>$ & $>$ &     & $>$ & $>$ \\
  $\$$                       & $<$ & $<$ & $<$ & $<$ & $<$ & $<$ & $<$ &     &     \\

  \bottomrule
  \end{tabular}
\end{table}

\end{document}
